\documentclass[journal]{IEEEtran}
\usepackage{color}
\usepackage{cite}
\usepackage{amsfonts}
\usepackage{amssymb}
\usepackage{amsmath}
\usepackage{amsthm}
\usepackage{stmaryrd}
\usepackage{url}
\usepackage{array}
\usepackage{booktabs}
\usepackage[utf8]{inputenc}
\usepackage[english]{babel}
\usepackage{algorithm}
\usepackage[noend]{algpseudocode}

\newtheorem{assumption}{Assumption}
\newtheorem{invariant}{Invariant}
\newtheorem{proposition}{Proposition}
\newtheorem{corollary}{Corollary}
\newtheorem{obligation}{Obligation}
\newcommand{\namedhead}{}
\newtheorem*{namedenv}{\namedhead}
\newenvironment{defD}[2]{\renewcommand{\namedhead}{Definition #1}%
  \begin{namedenv}[#2]}{\end{namedenv}}
\newenvironment{lawL}[2]{\renewcommand{\namedhead}{Law #1}%
  \begin{namedenv}[#2]}{\end{namedenv}}

\newcommand{\sem}[1]{\llbracket #1\rrbracket}
\newcommand{\then}{\mathbin{\ggg}}
\newcommand{\relay}{\mathsf{relay}}
\newcommand{\Int}{\mathsf{Intermediate}}
\newcommand{\WF}{\mathsf{WorldFacing}}
\newcommand{\Kl}{\mathbf{Kl}}
\newcommand{\Free}{\mathcal{F}}
\newcommand{\cl}{\diamond}
\newcommand{\Enc}[2]{\mathrm{E}^{#1}_{#2}}
\newcommand{\Dec}[2]{\mathrm{D}^{#1}_{#2}}
\newcommand{\view}{\mathrm{view}}
\newcommand{\nxt}{\mathrm{next}}
\newcommand{\frm}{\mathrm{from}}
\newcommand{\pad}{\mathrm{pad}}
\newcommand{\enc}{\mathrm{enc}}
\newcommand{\peel}{\mathrm{peel}}
\newcommand{\nf}{\mathrm{nf}}
\newcommand{\LoopOf}{\mathrm{Loop}}
\newcommand{\spec}{\mathrm{spec}}
\newcommand{\quota}{\mathrm{quota}}
\newcommand{\pos}{\mathrm{pos}}
\newcommand{\Args}{\mathrm{Args}}
\newcommand{\In}{\mathrm{In}}
\newcommand{\Out}{\mathrm{Out}}
\newcommand{\arr}{t^{\mathrm{arr}}}
\newcommand{\dep}{t^{\mathrm{dep}}}
\newcommand{\eqd}{\overset{d}{=}}
\DeclareRobustCommand{\code}[1]{{\ttfamily\renewcommand{\_}{\textunderscore\allowbreak}#1}}
\algnewcommand{\LineComment}[1]{\State \(\triangleright\) #1}

\author{Ryan~J.~Kung%
\thanks{Ryan J. Kung can be reached at ryankung@ieee.org.
Manuscript created September 24, 2026.  This document freezes the design
of Rings issue \#834 (\texttt{rings-node}, module \texttt{onion}) before
its Phase~2 code; it is a specification, not shipped behaviour.  The
segment relay count $s$ (\code{ONION\_SEGMENT\_RELAYS}) is left symbolic.}}
\title{\huge OnionMonad:\\
Onion Circuits as Client-Sealed Loops\\
of Registered Operation Symbols}
\date{}

\IEEEaftertitletext{%
\vspace{0.5\baselineskip}%
\noindent\hfill
\parbox{0.90\textwidth}{%
{\centering\bfseries Abstract\par}%
\vspace{0.4\baselineskip}%
\noindent
The Rings onion data plane interprets two layers, \emph{relay} and
\emph{exit}, and answers along the reversed path through relays that keep
per-circuit return state and forward one client-bound ciphertext
byte-identical.  We generalise the circuit to a term over a static
signature $\Sigma$ of operation symbols: nodes are partial
$\Sigma$-algebras that register symbols, never applications; a circuit is a
client-sealed loop $\cl\to g\to\cdots\to g\to\cl$ that evaluates one
application per hop and composes them in the Kleisli category of the hop
effect monad, with $\relay$ as the identity.  Every hop parses one uniform
layer; values travel in a fixed-width carry slot enciphered with AEZ, one
length-preserving layer per relay and one authenticated layer per
consumer; replies to world-facing sessions ride single-use reply blocks
under credit flow control.  We state ten laws, among them noninterference
over hop views, client path authority (which admits Arrow and Selective
composition and excludes monadic bind), byte unlinkability, and
single-use admission, list the residual leakage, and map each law to a
test obligation.\par\vspace{0.6\baselineskip}}}

\begin{document}
\maketitle

\section{Introduction}

Today's circuit is the term $\relay^{k}\then(\mathit{svc},\mathit{body})$
evaluated on the forward path and answered on its reverse
(\S\ref{sec:current}).  This paper states the general structure of which
it is the degenerate instance, as definitions D1--D9 and laws L1--L10 of
issue \#834, numbered identically.  Design lineage: Sphinx headers and
wide-block payloads~\cite{Danezis2009}, reply blocks~\cite{Chaum1981,
Danezis2003}, loops~\cite{Piotrowska2017}, AEZ~\cite{Hoang2015}, Kleisli
semantics~\cite{Moggi1991}, generic effects~\cite{Plotkin2003,
Plotkin2009}, and the Arrow/Selective hierarchy~\cite{Hughes2000,
Lindley2011,Mokhov2019}.

\section{Signature and Algebras}\label{sec:sig}

\begin{defD}{D1}{signature}
$\Sigma$ is a finite, statically defined table of symbols; each $f\in\Sigma$
has
\begin{multline*}
\spec(f)=(\mathit{name},\Args_f,\In_f,\Out_f,\\
W(f),L(f),\pos(f),\quota(f))
\end{multline*}
with $\mathit{name}$ identifying the backend (e.g.\ \code{jev.typesafe}),
$\Args_f$ client constants, $\In_f,\Out_f$ carried types, $W(f)\in\mathbb{N}$
the exact output width in bytes, $L(f)$ the exact latency class,
$\pos(f)\subseteq\{\Int,\WF\}$, and $\quota(f)$ the admission budget.
$W,L$ are protocol data.  The identity symbol is
\[
\relay:X\to X,\qquad W(\relay)=|X|,\qquad L(\relay)=0 .
\]
\end{defD}

\begin{defD}{D2}{registration}
Node $n$ is a partial $\Sigma$-algebra: a set $\Sigma_n\subseteq\Sigma$ and,
for each $f\in\Sigma_n$, a generic effect~\cite{Plotkin2003}
\[
\sem{f}_n:\Args_f\times\In_f\to M\,\Out_f ,
\]
$M$ the hop effect monad (failure, delay, world I/O).  $n$ publishes
$f\in\Sigma_n$, never an application.  $\relay$ is published in
\code{OnlineNodeDescriptor::capabilities}; every other $f$ in one signed
descriptor per symbol carrying $n$'s invocation epoch and $\quota(f)$.
\end{defD}

A node is a handler~\cite{Plotkin2009} for the operations it registers;
$\WF$ symbols (\code{tcp}, \code{https}) may hold a session, $\Int$ symbols
are one-shot.

\begin{defD}{D3}{application}
An application is $(f,\bar a)$, $\bar a\in\Args_f$, typed
$(f,\bar a):\In_f\rightsquigarrow\Out_f$; it is sealed in the layer of the
hop that evaluates it and visible to no other hop.
\end{defD}

\section{Pipelines}\label{sec:pipeline}

\begin{defD}{D4a}{pipeline}
Pipelines are the morphisms of the free category $\Free(\Sigma)$ on the
graph with types as vertices and applications as edges; composition is
\[
\frac{\alpha:X\rightsquigarrow Y\qquad \beta:Y\rightsquigarrow Z}
     {\alpha\then\beta:X\rightsquigarrow Z}.
\]
$P=(f_1,\bar a_1)\then\cdots\then(f_n,\bar a_n)$ is well formed iff
\begin{gather*}
n\ge1,\qquad\forall k<n:\ \Out_{f_k}=\In_{f_{k+1}},\\
\forall k<n:\ \Int\in\pos(f_k),\\
f_n\text{ is a }\WF\text{ session symbol}\ \Rightarrow\ n=1 .
\end{gather*}
Violations are rejected at construction.
\end{defD}

\begin{defD}{D4b}{Kleisli denotation}
For $\phi:X\to MY$, $\psi:Y\to MZ$ in $\Kl(M)$~\cite{Moggi1991},
$\phi\then\psi=\mu_Z\circ M\psi\circ\phi$.  For hop assignment
$h_1..h_n$ with $f_k\in\Sigma_{h_k}$,
\[
\sem{P}(v_0)=\bigl(\sem{f_1}_{h_1}(\bar a_1,-)\then\cdots\then
\sem{f_n}_{h_n}(\bar a_n,-)\bigr)(v_0),
\]
the unique functor $\Free(\Sigma)\to\Kl(M)$ extending the hop
interpretations.  The client consumes $\sem{P}(v_0)$.
\end{defD}

\begin{lawL}{L1}{identity}
$\sem{\relay}=\eta$, hence $\relay\then f=f=f\then\relay$ on values.  The
loop builder factors through $\nf$, which deletes $\relay$ applications:
$\LoopOf(P)=\LoopOf(\nf P)$, so $\LoopOf(\relay\then P)\eqd\LoopOf(P)$ and every
hop's view is unchanged.
\end{lawL}

\begin{lawL}{L2}{associativity}
$(f\then g)\then h=f\then(g\then h)$ on values; immediate from the monad
laws, and $\Free(\Sigma)$ has no bracketing.
\end{lawL}

Phase~1 admits only $\relay^{k}\then(f,\bar a)$ with $\WF\in\pos(f)$, which
is today's circuit (\S\ref{sec:current}).

\section{The Loop}\label{sec:loop}

\begin{defD}{D4c, D5}{loop, segment}
Let $\cl$ be the client, $g$ its entry guard, $s\ge1$ the protocol constant
\code{ONION\_SEGMENT\_RELAYS}.  Segment $k\in[0,n]$ is
$r_{k,1}..r_{k,s}$ from producer $p_k$ to consumer $c_k$,
\begin{gather*}
p_0=\cl,\quad p_k=h_k,\quad c_k=h_{k+1},\quad c_n=\cl,\\
\LoopOf(P):\ \cl\to r_{0,1..s}\to h_1\to r_{1,1..s}\to\cdots\\
\to h_n\to r_{n,1..s}\to\cl .
\end{gather*}
Positions are $i\in[1,N]$, $N=n+(n+1)s$; position $i$ evaluates $f_i$,
which is $\relay$ on every $r_{k,j}$.  $\LoopOf(\nf P)$ is the symbol word
$\relay^{s}f_1\relay^{s}\cdots f_n\relay^{s}$.
\end{defD}

\begin{lawL}{L7}{guard closure}
$r_{0,1}=r_{n,s}=g$, and the other $N-2$ positions carry pairwise distinct
DIDs, distinct from $g$.  $g$ is the only hop adjacent to $\cl$.
\end{lawL}

L7 requires $s\ge1$.  Route validation admits exactly the guard twice
(\code{has\_duplicate\_dids}), checks the loop shape against $s$, and
derives the hop bound from $N$.

\section{Uniform Layer and Timing}\label{sec:layer}

\begin{defD}{D6}{uniform layer}
Every position $i$ has one layer shape
\[
\lambda_i=(f_i,\bar a_i,\nxt_i,e_i,x_i,\nu_i,\kappa_i),
\]
epoch $e_i$, expiry $x_i$, replay nonce $\nu_i$, and carry keys
$\kappa_i=(k^{\mathrm{in}}_i;k^{\mathrm{out}}_{i,1..s+1})$ of constant
length; a relay has $\bar a_i=()$ and random $k^{\mathrm{out}}$.  The header
is Sphinx-shaped~\cite{Danezis2009}:
\[
\peel(\mathit{sk}_i,H_i)=(\lambda_i,H_{i+1}),\qquad |H_i|=|H| .
\]
The cell on edge $i\to\nxt_i$ is $(H_{i+1},y_i)\in\{0,1\}^{|H|}\times
\{0,1\}^{C}$ for every $i$.
\end{defD}

\begin{lawL}{L3}{width}
The carry slot has width $C$ on every edge.  With plaintext width
$C_0=C-\tau$, $W(f)<C_0$ for all $f\in\Sigma$ and
\[
\pad(b)=b\,\|\,\mathtt{80}\,\|\,\mathtt{00}^{C_0-|b|-1},\qquad
|\enc_f(w)|=W(f);
\]
an output with $|\enc_f(w)|>W(f)$ fails closed, never truncated.
\end{lawL}

\begin{lawL}{L4}{timing}
$\dep_i=\arr_i+L(f_i)$ exactly.  An evaluation unfinished at $\dep_i$ is
dropped, never released late.
\end{lawL}

\section{Carry}\label{sec:carry}

Write $\Enc{\tau}{k},\Dec{\tau}{k}$ for AEZ v5~\cite{AEZv5,Hoang2015}
encryption and decryption with empty nonce and associated data;
$|\Enc{\tau}{k}(m)|=|m|+\tau$, and $\Dec{\tau}{k}$ returns $\bot$ unless
the deciphered block ends in $\tau$ zero bytes.

\begin{defD}{D7}{carry}
For segment $k$ the client draws fresh keys
$k_{r_{k,1}},..,k_{r_{k,s}},k_{c_k}$ and places them
\[
\kappa_{p_k}^{\mathrm{out}}=(k_{r_{k,1}},..,k_{r_{k,s}},k_{c_k}),\quad
\kappa_{r_{k,j}}^{\mathrm{in}}=k_{r_{k,j}},\quad
\kappa_{c_k}^{\mathrm{in}}=k_{c_k}.
\]
With $b_k=\pad(\enc(v_k))$ and $\tau=16$:
\begin{gather*}
y_{k,0}=\Enc{0}{k_{r_{k,1}}}\circ\cdots\circ\Enc{0}{k_{r_{k,s}}}
\circ\Enc{\tau}{k_{c_k}}(b_k),\\
y_{k,j}=\Dec{0}{k_{r_{k,j}}}(y_{k,j-1}),\qquad
b_k=\Dec{\tau}{k_{c_k}}(y_{k,s}).
\end{gather*}
A consumer with $\In_f=()$ ignores the slot.
\end{defD}

Since $s$ is constant, $|\kappa_i|$ is constant and reveals no position.
By L9 a layer is admitted once, so every AEZ key enciphers exactly one
$C$-byte value, far below AEZ's birthday bound.

\begin{algorithm}[t]
\caption{$\mathrm{Hop}_i(\frm,H,y,\arr)$: one position, all symbols}
\label{alg:hop}
\begin{algorithmic}[1]
\State $(\lambda,H')\gets\peel(\mathit{sk}_i,H)$;\ \textbf{if} $\bot$
  \textbf{then} drop
\If{$\arr>x\ \vee\ e\ne\mathrm{epoch}_i\ \vee\ \nu\in R_i$} \State drop
  \Comment{L9} \EndIf
\State $R_i\gets R_i\cup\{(\nu,x)\}$
\If{$f=\relay$} \State $y'\gets\Dec{0}{k^{\mathrm{in}}}(y)$
\Else
  \State $b\gets\Dec{\tau}{k^{\mathrm{in}}}(y)$;\ \textbf{if} $b=\bot$
    \textbf{then} drop \Comment{L8}
  \State $w\gets\sem{f}(\bar a,\enc^{-1}\pad^{-1}(b))$ by $\arr+L(f)$
    \Comment{L4}
  \If{late $\vee\ |\enc_f(w)|>W(f)$} \State drop \Comment{L3} \EndIf
  \State $y'\gets\Enc{0}{k^{\mathrm{out}}_1}\circ\cdots\circ
    \Enc{0}{k^{\mathrm{out}}_s}\circ\Enc{\tau}{k^{\mathrm{out}}_{s+1}}
    (\pad\,\enc_f(w))$
\EndIf
\State release $(H',y')$ to $\nxt$ at $\arr+L(f)$
\end{algorithmic}
\end{algorithm}

\section{Reply Blocks}\label{sec:surb}

\begin{defD}{D8}{SURB}
A single-use reply block for a session symbol at $h$ is the client-sealed
loop suffix after $h$:
\[
\upsilon=(\nxt,H_\upsilon,k_{r_1},..,k_{r_s},k_{\cl},x_\upsilon).
\]
The suffix of the loop that opened or continued the session is itself a
SURB; a forward frame may carry $\sigma$ further SURBs in its carry value.
\end{defD}

\begin{invariant}[Credit]\label{inv:credit}
Let $Q_h$ be $h$'s unexpired SURBs, $u(t)$ the SURBs received and $r(t)$
the replies emitted by time $t$.  Then
\[
r(t)\le u(t),\qquad |Q_h|=0\Rightarrow\text{\code{tcp} stops reading its
socket.}
\]
Each reply consumes one $\upsilon$, produced as segment $n$ of
D7 with $p_n=h$; replenishment resumes reading.
\end{invariant}

\section{Security Laws}\label{sec:laws}

\begin{assumption}[Primitives]\label{as:prim}
Under a key it does not hold, a hop sees $\Enc{0}{k}$ as a uniformly random
permutation of $\{0,1\}^{C}$ and $\Enc{\tau}{k}$ as robust
AE~\cite{Hoang2015}; the header hides every layer but the hop's
own~\cite{Danezis2009}.  Network delays $\delta$ are fixed per edge.
\end{assumption}

\begin{defD}{D9}{view}
Over client coins $\rho$ (keys, nonces, header randomness) and delays
$\delta$,
\begin{multline*}
\view_i=(f_i,\bar a_i,v_{i-1},\frm_i,\nxt_i,e_i,x_i,\\
\mathit{bucket},\arr_i,\dep_i,H_i,y_{i-1},H_{i+1},y_i),
\end{multline*}
with $v_{i-1}=\bot$ when $f_i=\relay$ (a relay consumes no value).
\end{defD}

\begin{lawL}{L5}{noninterference}
For loops $\ell,\ell'$ let
\begin{gather*}
\pi_i(\ell)=(f_i,\bar a_i,v_{i-1},\frm_i,\nxt_i,e_i,x_i),\\
\sigma(\ell)=f_1\cdots f_N,\\
\pi_i(\ell)=\pi_i(\ell')\ \wedge\ \sigma(\ell)=\sigma(\ell')\ \Rightarrow\
\view_i(\ell)\eqd\view_i(\ell') .
\end{gather*}
Symbols may leak; $\bar a_j,v_j$ for $j\ne i$, other than the consumed
$v_{i-1}$, may not~\cite{Goguen1982}.
\end{lawL}

\begin{proposition}[L5 from D6, D7, L3, L4]\label{prop:ni}
Under Assumption~\ref{as:prim}, every loop built by
D4c--D7 satisfies L5.
\end{proposition}
\begin{proof}
Lengths: $|H_i|=|H|$, $|y_i|=C$ (D6, L3).  Times:
$\arr_i=\sum_{j<i}(L(f_j)+\delta_j)$ by L4, a function of $\sigma,\delta$.
Headers: $H_i,H_{i+1}$ reveal only $\lambda_i$, fixed by $\pi_i$ up to
fresh coins.  Carry: if $i=r_{k,j}$, then
$\Dec{0}{k_{r_{k,j}}}(y_{k,j-1})$ is an $\Enc{}{}$ output under
$k_{r_{k,j+1}}$ or $k_{c_k}$, a fresh key not held by $i$, hence uniform on
$\{0,1\}^C$ and independent of $v_k$.  If $i=c_k$, it learns $v_{i-1}$,
which $\pi_i$ fixes.  If $i=p_k$, its output is computed from $v_{i-1}$
and $\bar a_i$.  No component depends on $\bar a_j,v_j$ otherwise.
\end{proof}

\begin{lawL}{L6}{path authority}
$\nxt_i$ is a field of $\lambda_i$, fixed by $\cl$ at construction:
$\nxt_i$ is independent of every $v_j$.  No symbol chooses a hop.
\end{lawL}

\begin{proposition}[Admissible combinators]\label{prop:arrow}
L6 admits Arrow composition~\cite{Hughes2000} and Selective
branching~\cite{Mokhov2019} over client-sealed branches, and excludes
monadic bind.
\end{proposition}
\begin{proof}
$\then$ and $\mathit{select}:F(A+B)\to F(A\to B)\to F\,B$ fix the set of
possible effects before any value exists, so $\cl$ can seal every hop of
every branch in advance.  Bind $m\gg\!\!=k$, $k:A\to M\,B$, chooses the
continuation, and so its hops, from a run-time value: either a hop
constructs $\nxt$ (violating L6) or $\cl$ seals $|A|$ continuations.  This
is the arrows-meticulous, monads-promiscuous separation of
\cite{Lindley2011}.
\end{proof}

``OnionMonad'' names $M$, the effect inside one hop; across hops only the
static fragment of $\Kl(M)$ is exposed.

\begin{lawL}{L8}{unlinkability}
On consecutive edges $H_i\ne H_{i+1}$ and $y_{i-1}\ne y_i$ (one layer
peeled, one AEZ layer removed per hop).  If any $y_{k,j}$ is modified to
$\tilde y\ne y_{k,j}$,
\[
\Pr\bigl[\Dec{\tau}{k_{c_k}}(\tilde y_{k,s})\ne\bot\bigr]\le2^{-8\tau}
+\epsilon_{\mathrm{AEZ}},
\]
and a rejected value yields nothing; no bit of a tag survives.
\end{lawL}

Under Assumption~\ref{as:prim} consecutive $y$ are independent uniform
strings to any hop lacking both keys.  A modification propagates through
every later wide-block decipherment, so, unlike counter-mode relay
cryptography~\cite{Degabriele2018}, a flipped bit is never observed in
place; this is Sphinx's wide-block payload~\cite{Danezis2009,Anderson1996}
with AEZ in place of LIONESS.

\begin{lawL}{L9}{replay}
Every hop admits each $\nu$ at most once within $[\,\cdot\,,x]$; $R_i$
holds $(\nu,x)$ until $x$.  A SURB is a layer, hence single use.
\end{lawL}

\begin{lawL}{L10}{wire}
Phase~1 preserves today's bytes (golden tests); Phase~2 replaces them with
a new pinned golden set in one cutover, without versioning.
\end{lawL}

\section{Accepted Leakage}\label{sec:leak}

\begin{table}[t]
\caption{Residual leakage, stated and accepted}
\label{tab:leak}
\centering\small
\begin{tabular}{@{}>{\raggedright\arraybackslash}p{0.38\columnwidth}>{\raggedright\arraybackslash}p{0.55\columnwidth}@{}}
\toprule
leakage & bound \\
\midrule
symbols $\sigma$ are public; $\lambda_i$ names $f_i$ & L5 hypothesis \\
a $\WF$ hop knows it faces the world & no other hop learns whether symbols follow (L3, D6) \\
$g$ sees departure and return & loop round-trip time \\
latency class & L4 quantises timing; cannot remove it \\
hosted backend & sees $\bar a$ of every application of its symbol \\
$p_k$ holds $\kappa^{\mathrm{out}}_{p_k}$ & $p_k$ colluding with $r_{k,j}$ links two positions of one segment; no value leaks, and timing links them anyway \\
client-authored semantics ($\bar a_{i+1}$ depends on meaning of $v_i$) & out of protocol scope \\
\bottomrule
\end{tabular}
\end{table}

Table~\ref{tab:leak} is exhaustive with respect to L5: every other
component of $\view_i$ is determined by $\pi_i$ and $\sigma$.

\section{The Current Design}\label{sec:current}

\begin{proposition}[Backward linkability]\label{prop:current}
In \code{circuit/crypto.rs},
\code{encrypt\_client\_payload\_at\_sequence} seals the
exit's answer once, $c=\mathrm{AEAD}_{\mathit{pk}_\cl}(\cdot)$, and
in \code{circuit/reducer.rs}, \code{advance\_backward} re-emits it as
$\mathit{payload}=\mathit{frame.payload}$, rewriting only the edge
circuit id and the outer hop seal.  Every backward relay therefore sees the
same $c$: any two colluding relays, and the exit, link their positions by
equality, independently of route length.
\end{proposition}

\begin{table}[t]
\caption{Current circuit versus OnionMonad}
\label{tab:delta}
\centering\footnotesize
\setlength{\tabcolsep}{2pt}
\begin{tabular}{@{}>{\raggedright\arraybackslash}p{0.19\columnwidth}>{\raggedright\arraybackslash}p{0.39\columnwidth}>{\raggedright\arraybackslash}p{0.38\columnwidth}@{}}
\toprule
concern & current & here \\
\midrule
layers & \code{Relay}, \code{Exit} & one uniform $\lambda_i$ (D6) \\
computation & exit service only & any $f\in\Sigma_n$ (D2) \\
return path & forward path reversed & fresh relays to $g$ (L7) \\
relay state & \code{relay\_returns}, $\le1024$, TTL $120$\,s refreshed per cell & replay tags only (L9) \\
return bytes & one ciphertext, identical per edge & one AEZ layer per edge (L8) \\
replay & exits (\code{consume\_forward\_nonce}) & every hop (L9) \\
session replies & no credit bound & one per SURB (Inv.~\ref{inv:credit}) \\
\bottomrule
\end{tabular}
\end{table}

The stored return edge $(\mathit{id}',\nxt)\mapsto(\frm,\mathit{id},
\mathit{pk})$ is the only reason relays hold circuit state; the loop
removes it with \code{Backward}, \code{OnionBackwardFrame},
\code{MAX\_ONION\_RELAY\_CIRCUITS}, and \code{ONION\_RELAY\_RETURN\_TTL\_MS}.

\section{Security}\label{sec:security}

Table~\ref{tab:threats} maps each threat to the law that bounds it.

The adversary controls any set of nodes other than $\cl$, and delivery
order; it cannot break Assumption~\ref{as:prim}.  Dropping is always
possible and out of scope.

\begin{table}[t]
\caption{Threats and the law that bounds each}
\label{tab:threats}
\centering\small
\begin{tabular}{@{}>{\raggedright\arraybackslash}p{0.3\columnwidth}>{\raggedright\arraybackslash}p{0.63\columnwidth}@{}}
\toprule
threat & bound \\
\midrule
hop learns another's $\bar a_j$, $v_j$ & L5 (Prop.~\ref{prop:ni}) \\
position inference & constant $|H|$, $C$, $|\kappa|$, $s$ (D5--D7, L3) \\
value-dependent timing & exact release $\arr+L(f)$ (L4) \\
colluding relays match bytes & per-edge re-encipherment (L8); cf.\ Prop.~\ref{prop:current} \\
tagging & whole-value rejection, $2^{-8\tau}$ (L8) \\
hop-chosen routing & $\nxt_i\in\lambda_i$ (L6); bind excluded (Prop.~\ref{prop:arrow}) \\
first/last-hop correlation & one guard closes the loop (L7) \\
replay of layer or SURB & single admission per $\nu$ (L9) \\
reply amplification & $r(t)\le u(t)$ (Inv.~\ref{inv:credit}) \\
symbol abuse & per-symbol $\quota(f)$ (D1) \\
oversized result & fail closed at $W(f)$ (L3) \\
\bottomrule
\end{tabular}
\end{table}

\section{Verification}\label{sec:verification}

Every law is discharged against the pure builder and reducer, not a shadow
model (Table~\ref{tab:obligations}).

\begin{table}[t]
\caption{Obligations and the test that discharges each}
\label{tab:obligations}
\centering\small
\begin{tabular}{@{}>{\raggedright\arraybackslash}p{0.2\columnwidth}>{\raggedright\arraybackslash}p{0.73\columnwidth}@{}}
\toprule
obligation & test or property \\
\midrule
L10 & golden encodings of layers, wire messages, cells, and descriptor signing data, landed before any refactor; a new pinned set at the Phase~2 cutover \\
L1 & carry through $\relay$ is identity; $\LoopOf(\relay\then P)\eqd\LoopOf(P)$ \\
D4a & $\Out\ne\In$, $\WF$ before the last symbol, empty pipeline: rejected at construction \\
L3 & fixed slot width on every edge; output $>W(f)$ dropped \\
L4 & release at $\arr+L(f)$ and drop on overrun under an injected clock \\
L5 & property test on the pure builder: vary $\bar a_j,v_j$, $j\ne i$; $\view_i$ bytes, lengths, and release times equal \\
L7 & guard closure; any other duplicate DID rejected \\
L8 & per-hop byte inequality; one-bit flip rejects the whole value \\
L9 & replayed layer and reused SURB rejected \\
credit & a \code{tcp} session emits no reply without a SURB and resumes after replenishment \\
AEZ & official AEZ v5 vectors; \code{fuzz/} target for $\Dec{}{}\circ\Enc{}{}=\mathrm{id}$ and authenticator rejection \\
\bottomrule
\end{tabular}
\end{table}

\section{Phases}\label{sec:phases}

\begin{table}[t]
\caption{Cutover plan}
\label{tab:phases}
\centering\small
\begin{tabular}{@{}>{\raggedright\arraybackslash}p{0.14\columnwidth}>{\raggedright\arraybackslash}p{0.79\columnwidth}@{}}
\toprule
phase & content \\
\midrule
1 & golden bytes; $\Sigma$ table; pipeline term; layer build as a fold over $P$; reducer and handler as $\Sigma$-algebra; route by symbol sequence.  Zero wire change, no bump \\
2a & uniform layer, loop, stateless relays, AEZ carry, SURBs.  Breaking \\
2b & intermediate symbols: \code{Invoke} with release time, $\quota(f)$.  Same cutover; one version bump on its last PR \\
\bottomrule
\end{tabular}
\end{table}

Table~\ref{tab:phases} orders the cutover.  Open parameters: $s$ (proposed $1$), $\sigma$ and $C$, and the AEZ
implementation (candidate \code{zears}, unaudited, gated by
Table~\ref{tab:obligations}).  Non-goals: bind, Selective branches before
Phase~2b, session symbols in intermediate position, concrete symbols, and
verification of results.

\section{Related Work}\label{sec:related}

Tor~\cite{Dingledine2004} answers along the circuit it built, as Rings does
today; Degabriele and Stam~\cite{Degabriele2018} formalise the tagging that
its counter-mode relay cryptography admits.  Chaum's untraceable return
addresses~\cite{Chaum1981} and Mixminion's SURBs~\cite{Danezis2003} are the
source of D8; Sphinx~\cite{Danezis2009} gives the fixed-length header and a
LIONESS~\cite{Anderson1996} payload whose zero-block check L8 inherits.
Loopix~\cite{Piotrowska2017} routes traffic as loops for cover; here the
loop is the circuit.  AEZ~\cite{Hoang2015,AEZv5} supplies both the
length-preserving relay layer and robust AE in one primitive.  On the
semantic side, Moggi's Kleisli semantics~\cite{Moggi1991}, generic effects
and handlers~\cite{Plotkin2003,Plotkin2009}, and the
Arrow/Selective/Monad hierarchy~\cite{Hughes2000,Lindley2011,Mokhov2019}
fix which compositions a client can seal in advance.

\bibliographystyle{IEEEtran}
\bibliography{onion-monad}

\end{document}
